\documentclass[aspectratio=169,11pt]{beamer}
\usepackage{fontspec}
\usepackage{listings}
\usepackage{xcolor}
\usepackage{booktabs}
\usepackage{array}

% ================= Palette (identique au livre Tome 1) =================
\definecolor{PyBleu}{HTML}{1A4A7A}
\definecolor{PyBleuMoyen}{HTML}{2D7AB3}
\definecolor{PyBleuClair}{HTML}{EAF1F8}
\definecolor{PyJaune}{HTML}{FFD43B}
\definecolor{PyNuit}{HTML}{142840}
\definecolor{PyTexte}{HTML}{1F2A37}
\definecolor{PyGris}{HTML}{5A6B7E}
\definecolor{PyRouge}{HTML}{B3442D}
\definecolor{PyVert}{HTML}{2E7D5B}
\definecolor{CodeTexte}{HTML}{E8EEF6}

% ================= Code Python =================
\lstdefinestyle{py}{
  language=Python,
  keywordstyle=\color{PyJaune}\bfseries,
  stringstyle=\color{PyJaune},
  commentstyle=\color{PyGris}\itshape,
  showstringspaces=false,
  breaklines=true,
  backgroundcolor=\color{PyNuit},
  basicstyle=\ttfamily\footnotesize\color{CodeTexte},
  columns=fullflexible,
  keepspaces=true,
  xleftmargin=6pt,
  framexleftmargin=6pt,
}
\lstdefinestyle{pysmall}{
  language=Python,
  keywordstyle=\color{PyJaune}\bfseries,
  stringstyle=\color{PyJaune},
  commentstyle=\color{PyGris}\itshape,
  showstringspaces=false,
  breaklines=true,
  backgroundcolor=\color{PyNuit},
  basicstyle=\ttfamily\scriptsize\color{CodeTexte},
  columns=fullflexible,
  keepspaces=true,
  xleftmargin=5pt,
  framexleftmargin=5pt,
}

% ================= Titres de diapositives =================
\setbeamercolor{structure}{fg=PyBleu}
\setbeamerfont{frametitle}{series=\bfseries, size=\Large}
\setbeamercolor{frametitle}{bg=PyBleu, fg=white}
\setbeamertemplate{frametitle}{%
  \nointerlineskip
  \begin{beamercolorbox}[wd=\textwidth, ht=2.1ex, dp=0.9ex,
                          leftskip=0.35cm, rightskip=0.35cm]{frametitle}%
    \usebeamerfont{frametitle}\insertframetitle%
  \end{beamercolorbox}%
  \vspace*{-0.4pt}%
  {\color{PyJaune}\rule{\textwidth}{1.6pt}}%
}

% ================= Pied de page =================
\setbeamerfont{footline}{size=\tiny}
\setbeamercolor{footline}{bg=white, fg=PyGris}
\newcommand{\PiedDePage}{Python en 24 heures \textbar\ Séance 06 : Dict, set et fonctions}
\setbeamertemplate{footline}{%
  \leavevmode%
  \hbox{%
    \begin{beamercolorbox}[wd=\paperwidth, ht=2.2ex, dp=1.1ex,
                            leftskip=0.35cm, rightskip=0.35cm]{footline}%
      \usebeamerfont{footline}\PiedDePage
      \hfill\insertframenumber\,/\,\inserttotalframenumber%
    \end{beamercolorbox}}%
}

\setbeamertemplate{navigation symbols}{}
\setbeamertemplate{itemize item}{\color{PyBleuMoyen}$\blacktriangleright$}
\setbeamertemplate{itemize subitem}{\color{PyBleuMoyen}\textbullet}

% ================= Blocs pedagogiques =================
\setbeamercolor{block title}{bg=PyBleu, fg=white}
\setbeamercolor{block body}{bg=PyBleuClair, fg=PyTexte}
\setbeamercolor{block title alerted}{bg=PyRouge, fg=white}
\setbeamercolor{block body alerted}{bg=PyRouge!9, fg=PyTexte}
\setbeamercolor{block title example}{bg=PyJaune, fg=PyNuit}
\setbeamercolor{block body example}{bg=PyJaune!14, fg=PyTexte}

% ================= Metadonnees =================
\hypersetup{
  pdftitle={Python en 24 heures : Séance 06 : Dict, set et fonctions},
  pdfauthor={HOUNKPE Mawulikplimi Roland et DJABON Ounimborbitibou},
  pdfsubject={De l'algorithmique à la programmation orientée objet},
  pdfkeywords={Python, algorithmique, POO, formation, exercices corrigés}
}

\begin{document}

{%
\setbeamercolor{background canvas}{bg=PyNuit}
\begin{frame}[plain]
\vspace{10mm}
{\color{PyJaune}\small\bfseries FORMATION « PYTHON EN 24 HEURES » \quad\textbar\quad SÉANCE 06\par}
\vspace{4mm}
{\color{white}\Huge\bfseries Organiser ses données et son code : dict, set et fonctions\par}
\vspace{4mm}
{\color{PyBleuClair}\large Chapitres 11 et 12 du Tome 1\par}
\vspace{7mm}
{\color{PyJaune}\rule{34mm}{2.2pt}\par}
\vspace{7mm}
{\color{white}\normalsize 1 h de cours (présentation Beamer) \quad$\bullet$\quad 1 h de pratique (notebook Jupyter)\par}
\vspace{2.5mm}
{\color{PyBleuClair}\normalsize HOUNKPE Mawulikplimi Roland et DJABON Ounimborbitibou\par}
\end{frame}
}
\begin{frame}{Objectifs de la séance}
\begin{itemize}
  \item Créer et manipuler des dictionnaires : clés, valeurs, get et .items()
  \item Compter et regrouper des données avec un dictionnaire
  \item Utiliser les ensembles pour l'unicité et les opérations ensemblistes
  \item Écrire des fonctions avec paramètres, valeurs par défaut et return
  \item Documenter une fonction avec une docstring et choisir la bonne structure
\end{itemize}
\end{frame}
\begin{frame}{Rappel de la séance 05}
Les listes rangent par position. Quand la position n'a pas de sens, il faut une autre organisation : l'accès par clé.

\begin{itemize}
  \item Index, slices, méthodes append, insert, remove, pop
  \item Tuples : immuables, parfaits pour retourner plusieurs valeurs
  \item Compréhensions : simple, filtrée, imbriquée
  \item Piège : affecter une liste ne la copie pas
  \item enumerate pour parcourir avec la position
\end{itemize}
\end{frame}
\begin{frame}[fragile]{Le dictionnaire : accès par clé}

\begin{lstlisting}[style=py]
eleve = {"nom": "Awa", "age": 21}
print(eleve["nom"])            # Awa
eleve["filiere"] = "Maths"     # ajout
eleve["age"] = 22              # modification
print("nom" in eleve)          # True
print(eleve.get("tel"))        # None : pas d'erreur !
\end{lstlisting}
\vspace{0.5pt}

{\small\color{PyGris}\itshape Clé entre crochets : lecture et écriture. get(key) renvoie None au lieu de lever KeyError quand la clé manque.}

\end{frame}

\begin{frame}[fragile]{Parcourir un dictionnaire}
\begin{columns}[t]
\begin{column}{0.42\textwidth}
\begin{itemize}
  \setlength{\itemsep}{1.5pt}\setlength{\parskip}{0pt}
  \item `.keys()` : les clés (le parcours direct fait pareil)
  \item `.values()` : les valeurs seules
  \item `.items()` : les paires (clé, valeur), dépliées par l'unpacking
  \item Les dictionnaires conservent l'ordre d'insertion depuis Python 3.7
  \item Les clés sont uniques : réaffecter remplace la valeur
\end{itemize}
\end{column}
\hspace{0.02\textwidth}
\begin{column}{0.55\textwidth}
\begin{lstlisting}[style=pysmall]
stocks = {"riz": 50, "huile": 12}

for cle in stocks:
    print(cle, stocks[cle])

for cle, valeur in stocks.items():
    print(f"{cle:<8} : {valeur:>3} sacs")
\end{lstlisting}
\end{column}
\end{columns}
\end{frame}

\begin{frame}[fragile]{Compter avec un dictionnaire}

\begin{lstlisting}[style=py]
mots = ["riz", "mais", "riz", "yam", "mais", "riz"]
compteur = {}
for mot in mots:
    if mot in compteur:
        compteur[mot] += 1
    else:
        compteur[mot] = 1
print(compteur)
\end{lstlisting}
\vspace{0.5pt}

{\small\color{PyGris}\itshape Le schéma de comptage : si la clé existe, incrémenter ; sinon, créer à 1. La méthode setdefault ou .get(k, 0) + 1 condense ce code.}

\end{frame}

\begin{frame}[fragile]{L'ensemble : unicité garantie}
\begin{columns}[t]
\begin{column}{0.42\textwidth}
\begin{itemize}
  \setlength{\itemsep}{1.5pt}\setlength{\parskip}{0pt}
  \item Création : `set([1, 2, 2, 3])` donne \{1, 2, 3\}
  \item Aucun doublon, aucun ordre garanti
  \item `in` est \textbf{très rapide}, même sur de gros ensembles
  \item union `|`, intersection `\&`, différence `-`
  \item Usage typique : dédoublonner une liste
\end{itemize}
\end{column}
\hspace{0.02\textwidth}
\begin{column}{0.55\textwidth}
\begin{lstlisting}[style=pysmall]
inscrits = ["awa", "kofi", "awa", "sena"]
uniques = set(inscrits)
print(uniques)

a = {1, 2, 3}
b = {2, 3, 4}
print(a | b)   # {1, 2, 3, 4}
print(a & b)   # {2, 3}
print(a - b)   # {1}
\end{lstlisting}
\end{column}
\end{columns}
\end{frame}

\begin{frame}{Quelle structure choisir ?}
\renewcommand{\arraystretch}{1.3}\small
\begin{center}
\begin{tabular}{lll}
\toprule
\textbf{\color{PyBleu}Structure} & \textbf{\color{PyBleu}Quand l'utiliser} & \textbf{\color{PyBleu}Exemple} \\
\midrule
list & série ordonnée, modifiable & file d'attente, notes \\
tuple & groupe figé de valeurs & coordonnées, retour multiple \\
dict & accès par clé signifiante & fiche élève, inventaire \\
set & unicité et appartenance rapide & dédoublonner, tags \\
\bottomrule
\end{tabular}
\end{center}

\vspace{4pt}

{\small\color{PyGris}Règle pratique : la position compte ? list. Le nom compte ? dict. Le doublon interdit ? set.}
\end{frame}

\begin{frame}[fragile]{Définir une fonction}

\begin{lstlisting}[style=py]
def aire_rectangle(longueur, largeur):
    """Retourne l'aire d'un rectangle."""
    return longueur * largeur

resultat = aire_rectangle(8, 5)
print(resultat)   # 40
\end{lstlisting}
\vspace{0.5pt}

{\small\color{PyGris}\itshape def lance la définition, la docstring documente, return rend la valeur. Une fonction sans return rend None.}

\end{frame}

\begin{frame}[fragile]{Paramètres : ordre, défauts, mots clés}
\begin{columns}[t]
\begin{column}{0.42\textwidth}
\begin{itemize}
  \setlength{\itemsep}{1.5pt}\setlength{\parskip}{0pt}
  \item L'ordre des arguments positionnels compte
  \item Un paramètre avec défaut devient \textbf{facultatif}
  \item Les arguments par mots clés rendent l'appel explicite
  \item Les valeurs par défaut sont évaluées \textbf{une seule fois}
  \item Jamais de liste mutable en défaut (piège du Tome 1)
\end{itemize}
\end{column}
\hspace{0.02\textwidth}
\begin{column}{0.55\textwidth}
\begin{lstlisting}[style=pysmall]
def prix_ttc(hors_taxe, taux=0.18, remise=0):
    net = hors_taxe * (1 + taux)
    return round(net * (1 - remise), 2)

print(prix_ttc(100))
print(prix_ttc(100, taux=0.075))
print(prix_ttc(100, remise=0.1, taux=0.075))
\end{lstlisting}
\end{column}
\end{columns}
\end{frame}

\begin{frame}[fragile]{Portée : locale par défaut}

\begin{lstlisting}[style=py]
total_global = 1000     # variable globale

def ajouter(montant):
    part_locale = montant      # locale à la fonction
    return part_locale

print(ajouter(50))
print(total_global)
# print(part_locale)  -> NameError : hors de la fonction
\end{lstlisting}
\vspace{0.5pt}

{\small\color{PyGris}\itshape Une variable créée dans une fonction disparaît à la sortie. Une fonction lit les globales mais doit les éviter : passez tout en paramètres.}

\end{frame}

\begin{frame}[fragile]{Retourner plusieurs valeurs}

\begin{lstlisting}[style=py]
def min_max_moy(notes):
    """Retourne (minimum, maximum, moyenne) d'une série."""
    return (min(notes), max(notes),
            round(sum(notes) / len(notes), 2))

mini, maxi, moy = min_max_moy([12, 15, 9, 18])
print(mini, maxi, moy)
\end{lstlisting}
\vspace{0.5pt}

{\small\color{PyGris}\itshape Le tuple retourné se déplie en trois variables à l'appel : c'est le idiome standard du calcul de statistiques.}

\end{frame}

\begin{frame}{Attention}
\begin{alertblock}{KeyError : la clé fantôme}
`eleve["tel"]` lève `KeyError` si la clé n'existe pas. Deux parades : `.get("tel")` qui renvoie None, ou le test `"tel" in eleve` avant l'accès. Choisissez selon l'intention : get quand l'absence est normale, in quand l'absence est une anomalie à signaler.
\end{alertblock}
\end{frame}

\begin{frame}{À retenir}
\begin{block}{L'essentiel de la séance 06}
Le dictionnaire relie une clé à une valeur et se parcourt avec .items() ; l'ensemble garantit l'unicité et teste l'appartenance en un éclair. La structure se choisit selon la question posée : position, nom ou unicité. La fonction encapsule un calcul derrière un nom clair, documenté par une docstring, avec des paramètres explicites : c'est la brique de base de tout gros programme.
\end{block}
\end{frame}

\begin{frame}{Travaux pratiques (1 h) : Dictionnaires, ensembles et fonctions}
Cinq exercices pour organiser les données puis le code : fiche contact, comptage, validation, statistiques et un mini quiz complet. Les fonctions des exercices 3 à 5 doivent avoir une docstring d'une ligne.

\begin{itemize}
  \setlength{\itemsep}{1.5pt}\setlength{\parskip}{0pt}
  \item \textbf{Exercice 1 (Basique)} : Fiche contact
  \item \textbf{Exercice 2 (Basique)} : Compteur d'occurrences
  \item \textbf{Exercice 3 (Intermédiaire)} : Validation d'adresse e-mail
  \item \textbf{Exercice 4 (Intermédiaire)} : Statistiques de notes en dictionnaire
  \item \textbf{Exercice 5 (Défi)} : Mini quiz en fonctions
\end{itemize}
\vspace{4pt}

{\small\color{PyGris}Ouvrez le notebook « Seance\_06\_dict\_set\_fonctions.ipynb » pour commencer.}
\end{frame}
\begin{frame}{Pour la prochaine séance}
\begin{exampleblock}{Conseil de travail}
Séance 07 : fonctions avancées (closures et décorateurs), puis chaînes et fichiers. Relisez la portée locale et globale de cette séance : la closure est exactement une fonction qui garde en mémoire la portée où elle est née.
\end{exampleblock}
\end{frame}

\end{document}
